%-------------------------------------------------------------------------------
%	SECTION TITLE
%-------------------------------------------------------------------------------
\cvsection{Ongoing Works/Submissions}

%-------------------------------------------------------------------------------
%	CONTENT
%-------------------------------------------------------------------------------
\begin{cventries}

%---------------------------------------------------------

\cventry
    {Petros Drineas, Rafail Ostrovsky, Rohit Nema, Vassilis Zikas} % Role
    {Trustworthy Reputation Games and Applications to Proof-of-Reputation Blockchains} % Event
    {} % Location
    {submitted to EC'25} % Date(s)
    {
      \begin{cvitems} % Description(s)
        \item Building a Game-theoretic (rational, incentive-compatible) mechanism to formalize ``reputation'' in peer-to-peer networks. 
        \item Introduced a new model (class of games) called \emph{trustworthy reputation games} such that it is best response to play a prescribed (truthful) strategy according to one's beliefs
        \item Uses a PageRank in a clever way to model trustworthiness and utility of nodes in the network
        \item Applications include a new consensus mechanism for blockchain, tackling misinformation and improving reviews on platforms such as Amazon.
      \end{cvitems}
    }

%---------------------------------------------------------
  \cventry
    {Dan Boneh, Joshua Hsieh, Rohit Nema} % Role
    {Efficiently Auditable Key Transparency Using Algebraic Techniques} % Event
    {Stanford, CA} % Location
    {ongoing} % Date(s)
    {
      \begin{cvitems} % Description(s)
        \item Improving the download bandwidth of auditors in an end-to-end encrypted key transparency mechanism by several orders of magnitude while preserving other metrics for a practical real-world deployment.
        \item Key Transparency is currently deployed by huge platforms such as Apple's iMessage and Meta's WhatsApp and thus, any improvement can have major impact.
      \end{cvitems}
    }

%---------------------------------------------------------

\cventry
    {Dan Boneh, Rohit Nema} % Role
    {Expressive Set Accumulators} % Event
    {Stanford, CA} % Location
    {ongoing} % Date(s)
    {
      \begin{cvitems} % Description(s)
        \item{Building Set Accumulators which support \emph{expressive} operations such as union, intersection, min, max and range queries while keeping the CRS small (preferably, sublinear)}
        \item Instantiating using Multilinear Polynomial Commitments
      \end{cvitems}
    }

%---------------------------------------------------------

\cventry
    {Sam Dittmer, Rafail Ostrovsky, Rohit Nema} % Role
    {Maliciously Secure Secret-Shared Shuffling with Optimal Communication} % Event
    {} % Location
    {ongoing} % Date(s)
    {
      \begin{cvitems} % Description(s)
      	\item Developing a cryptographic protocol (2PC) to randomly shuffle a vector inputted as authenticated secret shares (SPDZ-like).
        \item Improving on the state-of-the-art constructions of the semi-honest \href{https://eprint.iacr.org/2019/1340}{[CGP'19]} and its maliciously secure version \href{https://eprint.iacr.org/2023/1794}{[SYBDC'23]} using linearly homomorphic encryption. 
        \item Optimal communication and computation complexity with a constant number of rounds.
      \end{cvitems}
    }

%---------------------------------------------------------
\end{cventries}
